% Materiał do adaptacji do rozdziału 8. Źródło: pomiary własne,
% results/reports/campaign_2026-09-20; identyfikacja źródeł: source_manifest.csv.
% Tabele nie wymagają dodatkowych pakietów. Rysunki: pakiet graphicx,
% skopiować PDF do katalogu grafik pracy i dostosować ścieżki includegraphics.
\section{Weryfikacja i wyniki pomiarów stanowiskowych}

Badania objęły 19 realizacji macierzy scenariuszy dla sześciu ziaren generatora,
łącznie 323 zaplanowane epizody. Osobno oceniono test funkcjonalny smoke\_03
oraz sześć kompletnych sesji uczenia profilu. Wyniki w tej części pochodzą
z pomiarów własnych. Zestawienie danych źródłowych i sum kontrolnych zapisano
w pliku \texttt{source\_manifest.csv}, a wyniki pojedynczych epizodów
w pliku \texttt{episodes.csv} towarzyszącym analizie.

Znaczniki początku i końca przesyłane na magistrali umożliwiły powiązanie
planu generatora z czasem urządzenia. W 18 z 19 macierzy uzyskano zgodność
kolejności, identyfikatorów, długości i danych każdej potwierdzonej ramki.
W pozostałej próbie nie odnaleziono dwóch ramek ruchu referencyjnego.
W śladach odbiornika potwierdzono wszystkie 11\,110 ramek przypisanych
do wstrzykniętych anomalii i potwierdzonych przez adapter.

\begin{table}[htbp]
\centering
\small
\caption{Liczba realizacji z nowym alarmem w epizodzie. Źródło: pomiary własne.
Brak ponownego alarmu DBC może oznaczać utrzymanie stanu naruszenia.}
\label{tab:wyniki-macierz-alarmy}
\setlength{\tabcolsep}{3pt}
\begin{tabular}{p{0.07\linewidth}p{0.47\linewidth}p{0.20\linewidth}p{0.16\linewidth}}
\hline
Epizod & Wariant & Reguła podstawowa & Dowolna reguła \\
\hline
E01 & Wzrost 1,25$\times$ & 0/19 & 19/19 \\
E02 & Wzrost 2$\times$ & 19/19 & 19/19 \\
E03 & Wzrost 4$\times$ & 19/19 & 19/19 \\
E04 & Spadek do 0,8$\times$ & 0/19 & 19/19 \\
E05 & Spadek do 0,5$\times$ & 13/19 & 19/19 \\
E06 & Spadek do 0,25$\times$ & 19/19 & 19/19 \\
E07 & Zanik 1 s & 19/19 & 19/19 \\
E08 & Zanik 5 s & 19/19 & 19/19 \\
E09 & Burst 4 ramki / 20 ms & 8/19 & 19/19 \\
E10 & Burst 20 ramek / 10 ms & 19/19 & 19/19 \\
E11 & Nieznane ID, 1 ramka & 19/19 & 19/19 \\
E12 & Nieznane ID, 8 ramek & 19/19 & 19/19 \\
E13 & DLC krótsze & 19/19 & 19/19 \\
E14 & DLC dłuższe & 16/19 & 16/19 \\
E15 & Wartość 101 (maks. 100) & 19/19 & 19/19 \\
E16 & Wartość -101 (min. -100) & 16/19 & 16/19 \\
E17 & Wartość 110 (maks. 100) & 13/19 & 13/19 \\
\hline
\end{tabular}
\end{table}

Nowy alarm co najmniej jednej metody wystąpił w 311 z 323 epizodów.
Wskaźnik ten opisuje rejestrację nowych zdarzeń, a nie czułość klasyfikacji
pojedynczych ramek. Reguły zapisują wejście w stan naruszenia i nie powtarzają
alarmu przed jego skasowaniem poprawną obserwacją. Dwanaście epizodów bez
nowego alarmu DBC było związanych z utrzymaniem takiego stanu pomiędzy
wstrzyknięciami. Dla wiadomości o okresie 10~s odstęp 8~s między początkami
epizodów nie gwarantował wystąpienia poprawnej ramki przywracającej stan normalny.

Detektor zaniku zareagował we wszystkich 38 próbach ciszy o długości 1~s
lub 5~s. Reguły odstępu mogły wykryć wydłużenie przerwy dopiero po powrocie
wiadomości. Reguła częstotliwości wykryła silniejsze zmiany w 57 z 57 prób
wzrostu dwukrotnego, czterokrotnego lub spadku do jednej czwartej wartości
odniesienia. Dla spadku do połowy częstotliwości wynik wynosił 13 z 19.
Ostra nierówność względem progu bliskiego pięciu ramkom powodowała,
że niewielka różnica wyuczonego okresu zmieniała decyzję na granicy zakresu.

W granicznym burście czterech ramek co 20~ms reguła burst zareagowała
w 8 z 19 prób, mimo potwierdzenia odbioru wszystkich wstrzykniętych ramek.
Rzeczywiste odstępy dochodziły do około 30~ms, przekraczając próg około 25~ms.
W gęstym burście co 10~ms uzyskano reakcję w 19 z 19 prób, lecz odebrano
356 z 380 zaplanowanych ramek, ponieważ harmonogram generatora pominął
24 terminy. W ocenie intensywności wykorzystano zatem rzeczywisty zapis odbioru.

% Proponowane rysunki: 01_macierz_metod.pdf, 02_realizacja_burstu.pdf,
% 03_alarmy_tla.pdf, 04_granica_czestotliwosci.pdf, 05_jitter_z_pomiaru.pdf.
% Dokładne wskaźniki alarmów tła i parametry wydajności: RAPORT_ANALIZY.md.
% Nie utożsamiać nowych alarmów/epizod z recall ramek ani alarmów/min z FPR.
